\section{Introducción}

El procesamiento del lenguaje natural (NLP, por sus siglas en inglés) ha experimentado una transformación radical desde la introducción de la arquitectura Transformer. Estos modelos, basados en mecanismos de atención, han demostrado un desempeño excepcional en una amplia variedad de tareas, superando las limitaciones de las arquitecturas recurrentes y convolucionales previas.

La presente práctica tiene como objetivo principal explorar en profundidad la potencia y las limitaciones de los modelos basados en Transformers aplicados a tareas complejas de NLP. A través de un enfoque práctico y experimental, se busca no solo implementar soluciones funcionales, sino también comprender los fundamentos de su arquitectura, el proceso de entrenamiento y la interpretabilidad de los resultados obtenidos.

Para ello, el trabajo se estructura en torno a la implementación y el ajuste fino (\textit{fine-tuning}) de al menos dos arquitecturas de Transformers distintas. El núcleo del estudio consiste en realizar una evaluación comparativa y un análisis crítico que abarque desde la optimización de hiperparámetros (como tasas de aprendizaje y estrategias de regularización) hasta el uso de métricas avanzadas de evaluación, incluyendo el uso de modelos de lenguaje de gran escala (LLMs) como jueces de calidad.

Finalmente, se abordará la importancia de la explicabilidad en estos modelos mediante técnicas de visualización de mapas de atención y métodos como \textit{Integrated Gradients} o \textit{SHAP}, permitiendo una comprensión más profunda de cómo los modelos procesan y generan información lingüística.
